\documentclass[11pt]{article}

\usepackage{nmrdoc}
\usepackage{siunitx}
\usepackage{bm}

\sisetup{per-mode=symbol}
\DeclareSIUnit{\angstrom}{\text{\AA}}
\DeclareSIUnit{\bohr}{a_0}
\DeclareSIUnit{\hartree}{Ha}
\DeclareSIUnit{\elementarycharge}{e}

\newcommand{\code}[1]{\texttt{#1}}
\setlength{\parindent}{0pt}
\setlength{\parskip}{0.42em}
\emergencystretch=2em

\begin{document}
\NmrTitle{Numerical Addendum: EEQ Calculator}

The EEQ calculator turns one conformation into two scalar vectors. The
supporting work is the smooth coordination-number count and the construction of
the dense Coulomb matrix. The numerical centre is the \(N\times N\) Cholesky
factorization of that matrix and the bordered-system reduction that enforces
the total charge without factoring the indefinite augmented matrix. It is the
only \(O(N^3)\) dense step in this calculator; the coordination count and
matrix fill are \(O(N^2)\), and the two solves reuse the same triangular
factor. The listings below are distilled from the source --- loops and
bookkeeping trimmed --- but the constants, guards, inequalities, and array
shapes are the implementation's.

\section{Smooth coordination counts}

The code first moves positions from \(\si{\angstrom}\) into \(\si{\bohr}\) and
looks up one D4 parameter row per atom. In the coordination-number pass,
\code{cn(i)} is the accumulated smooth neighbour count for atom \(i\). Each
unordered atom pair \(i,j\) is visited once. If its squared distance is greater
than \code{cn\_cutoff\_bohr\_sq}, the pair is absent from both counts. Otherwise
the code evaluates one number,
\[
  c_{ij}=
  \frac12\,\operatorname{erfc}
  \left[
    k_{\mathrm{CN}}
    \left(
      \frac{R_{ij}}{r_{\mathrm{cov},i}+r_{\mathrm{cov},j}}-1
    \right)
  \right],
\]
and adds that same \(c_{ij}\) to \code{cn(i)} and \code{cn(j)}.
\(R_{ij}\) is the pair distance in \(\si{\bohr}\),
\(r_{\mathrm{cov}}\) is the covalent radius from \code{params}, and
\(k_{\mathrm{CN}}=\code{eeq\_cn\_steepness}=7.5\). The comparison is strict, so
a pair landing exactly at the cutoff is still counted.

The complementary error function replaces a discontinuous integer neighbour
count by a smooth transition. At
\(R_{ij}=r_{\mathrm{cov},i}+r_{\mathrm{cov},j}\) the argument is zero, so the
pair contributes \(1/2\). Shorter pairs contribute more than \(1/2\) and
approach one as the argument becomes negative; longer pairs contribute less
than \(1/2\) and approach zero. The shape is differentiable in \(R_{ij}\) for
counted pairs. The configured \(\SI{25.0}{\angstrom}\) cutoff remains a hard
cutoff: a pair beyond it contributes exactly zero, with no reweighting of the
remaining pairs.

\begin{codebox}
double cn_cutoff_bohr = cn_cutoff * BOHR_PER_ANGSTROM;
double cn_cutoff_bohr_sq = cn_cutoff_bohr * cn_cutoff_bohr;

Eigen::VectorXd cn = Eigen::VectorXd::Zero(N);
for (size_t i = 0; i < N; ++i) {
    for (size_t j = i + 1; j < N; ++j) {
        double dx = pos(i,0) - pos(j,0);
        double dy = pos(i,1) - pos(j,1);
        double dz = pos(i,2) - pos(j,2);
        double dist_bohr_sq = dx*dx + dy*dy + dz*dz;
        if (dist_bohr_sq > cn_cutoff_bohr_sq) continue;

        double dist_bohr = std::sqrt(dist_bohr_sq);
        double rcov_sum = params[i].rcov + params[j].rcov;
        double arg = cn_steepness * (dist_bohr / rcov_sum - 1.0);
        double count = 0.5 * std::erfc(arg);

        cn(i) += count;
        cn(j) += count;
    }
}

for (size_t i = 0; i < N; ++i)
    chi_eff(i) = params[i].chi
               + params[i].kappa * std::sqrt(cn(i) + 1e-14);
\end{codebox}

The second loop turns \code{cn(i)} into the effective electronegativity
\(\chi_i^{\mathrm{eff}}\) held in \code{chi\_eff(i)}: a base electronegativity
\(\chi_i\) (\code{params[i].chi}) shifted by the coordination through
\(\kappa_i\) (\code{params[i].kappa}), the per-element sensitivity of
electronegativity to neighbour count. The \(10^{-14}\) term is the source guard
carried from dftd4 arithmetic: an isolated atom still takes a square root of a
positive number. The guard changes the zero-coordination value by
\(\kappa_i\,10^{-7}\) in \(\si{\hartree}\).

\section{Dense Coulomb matrix}

\code{A} is a dense real \(N\times N\) matrix. Filling it is an \(O(N^2)\)
all-pairs pass. Its diagonal stores a local
quadratic charge cost,
\[
  A_{ii} = \eta_i + \frac{\sqrt{2/\pi}}{r_{\mathrm g,i}},
\]
where \(\eta_i\) is \code{params[i].gam} and \(r_{\mathrm g,i}\) is
\code{params[i].rad}. The first term is the chemical hardness. The second term
is the Gaussian self-Coulomb term. Both are positive for every project-local
protein element row and for the fallback row.

For \(i\ne j\), \code{A(i,j)} is the Ohno--Klopman pair kernel
\[
  A_{ij} =
  \gamma_{ij}(R_{ij}) =
  \frac{1}{\sqrt{R_{ij}^2+1/(\eta_i\eta_j)}} .
\]
The same value is written into \code{A(j,i)}, so the stored matrix is symmetric.
At zero separation the off-diagonal value is finite,
\(\gamma_{ij}(0)=\sqrt{\eta_i\eta_j}\); at long range it approaches the
atomic-unit Coulomb factor \(1/R_{ij}\).

\begin{codebox}
Eigen::MatrixXd A(N, N);
for (size_t i = 0; i < N; ++i) {
    A(i, i) = params[i].gam + SQRT_2_OVER_PI / params[i].rad;

    for (size_t j = i + 1; j < N; ++j) {
        double dx = pos(i,0) - pos(j,0);
        double dy = pos(i,1) - pos(j,1);
        double dz = pos(i,2) - pos(j,2);
        double dist_bohr_sq = dx*dx + dy*dy + dz*dz;

        double inv = 1.0 / (params[i].gam * params[j].gam);
        double gamma = 1.0 / std::sqrt(dist_bohr_sq + inv);
        A(i, j) = gamma;
        A(j, i) = gamma;
    }
}
\end{codebox}

The self term is the diagonal contribution from a charge's finite-width Coulomb
energy. Without it, identical atoms at zero separation would have
\(A_{ii}=\eta_i\) and \(A_{ij}=\eta_i\), so the diagonal would only tie one
contact off-diagonal. Adding \(\sqrt{2/\pi}/r_{\mathrm g,i}\) raises the
Gershgorin row centre above that contact value. For the H, C, N, O, S, and
fallback rows in \code{PhysicalConstants.h}, \(A_{ii}\) is larger than every
zero-distance off-diagonal involving atom \(i\). Gershgorin's theorem locates
every eigenvalue inside a disc --- one per row, centred on that row's diagonal
entry \(A_{ii}\) and reaching out by the sum of the row's off-diagonal
magnitudes. When the full row also
satisfies \(A_{ii}>\sum_{j\ne i}|A_{ij}|\), Gershgorin discs lie on the
positive real axis; because \(A\) is symmetric, its eigenvalues are then
positive, and the matrix is symmetric positive-definite. The implementation
does not separately test the row sums. It uses the factorization call below as
the SPD gate.

\section{Cholesky factorization}

Eigen's \code{LLT} call factors \(\bm A\) as
\(\bm A=\bm L\bm L^{\mathsf T}\), with \(\bm L\) lower triangular. The
factorization is the \(O(N^3)\) step in this calculator. Once \(\bm L\) exists,
each solve with a right-hand side \(\bm b\) is two triangular substitutions:
first solve \(\bm L\bm y=\bm b\) by forward substitution, then solve
\(\bm L^{\mathsf T}\bm x=\bm y\) by back-substitution. That returns
\(\bm x=\bm A^{-1}\bm b\) to the accuracy of the floating-point factor and
substitutions.

The code uses Cholesky rather than a general linear solve because this
formulation is intended to supply a symmetric positive-definite block. In exact
arithmetic, an SPD matrix gives positive pivots and no row exchanges are needed.
In this implementation, \code{chol.info()} is the numerical gate: if Eigen
cannot complete the positive-definite factorization on the supplied \code{A},
the calculator returns no result.

\begin{codebox}
Eigen::LLT<Eigen::MatrixXd> chol(A);
if (chol.info() != Eigen::Success) {
    OperationLog::Error("EeqResult",
        "Cholesky decomposition failed (N=" + std::to_string(N) +
        ") -- Coulomb matrix not positive definite");
    return nullptr;
}

Eigen::VectorXd A_inv_chi = chol.solve(chi_eff);
Eigen::VectorXd ones = Eigen::VectorXd::Ones(N);
Eigen::VectorXd A_inv_ones = chol.solve(ones);
\end{codebox}

The same factorization serves both right-hand sides. \code{A\_inv\_chi} is
\(\bm u\), the solution of \(\bm A\bm u=\bm\chi^{\mathrm{eff}}\).
\code{A\_inv\_ones} is \(\bm v\), the solution of \(\bm A\bm v=\bm 1\). The
cost of forming \(\bm L\) is paid once; the second right-hand side reuses the
same triangular factors.

\section{Charge constraint}

The constrained first-order system is bordered:
\[
  \begin{bmatrix}
    \bm A & \bm 1\\
    \bm 1^{\mathsf T} & 0
  \end{bmatrix}
  \begin{bmatrix}\bm q\\ \lambda\end{bmatrix}
  =
  \begin{bmatrix}-\bm\chi^{\mathrm{eff}}\\ Q\end{bmatrix}.
\]
\(Q\) is \code{Q\_total}, defaulting to
\(\SI{0}{\elementarycharge}\), and \(\lambda\) is the scalar multiplier for the
charge sum. The augmented matrix is indefinite because its lower-right block is
zero and the constraint row has no quadratic charge cost.
Positive-definiteness is a bowl, curving upward in every direction; the zero
lower-right block gives the multiplier \(\lambda\) a direction with no curvature
of its own, tipping the bordered system into a saddle --- up some ways, down
others. That saddle is what \emph{indefinite} names, and it is why Cholesky,
which needs the bowl, cannot factor it. The code avoids
factoring that matrix. It eliminates against the SPD block \(\bm A\).

From the first block row,
\[
  \bm q = -\bm A^{-1}\bm\chi^{\mathrm{eff}}
          -\lambda\bm A^{-1}\bm 1
        = -(\bm u+\lambda\bm v).
\]
Applying the second block row \(\bm 1^{\mathsf T}\bm q=Q\) gives
\[
  \lambda =
  -\frac{Q+\bm 1^{\mathsf T}\bm u}{\bm 1^{\mathsf T}\bm v}.
\]
Thus the two Cholesky solves build the two vectors used by the constraint:
\(\bm u\) is the unconstrained response to electronegativity, and \(\bm v\) is
the response to the all-ones vector in the first block row. The multiplier
chooses how much of \(\bm v\) must be mixed in so that the entries of \(\bm q\)
sum to \(Q\).

\begin{codebox}
double denom = ones.dot(A_inv_ones);        // 1^T v
if (std::abs(denom) < 1e-30) {
    OperationLog::Error("EeqResult",
        "charge constraint denominator is zero");
    return nullptr;
}

double lambda =
    -(Q_total + ones.dot(A_inv_chi)) / denom;

Eigen::VectorXd charges =
    -(A_inv_chi + lambda * A_inv_ones);

double charge_residual = charges.sum() - Q_total;
charges.array() -= charge_residual / static_cast<double>(N);
\end{codebox}

\code{denom} is \(\bm 1^{\mathsf T}\bm v\). If it is smaller than \(10^{-30}\)
in magnitude, the scalar constraint solve is treated as unusable and the
calculator returns no result; equality with \(10^{-30}\) proceeds. Otherwise the
algebraic reduction gives
\code{charges}. In exact arithmetic, those charges satisfy both the stationarity
equations and the charge constraint. In floating point, the two triangular
solves can leave a nonzero scalar charge residual; the code does not estimate
\(\kappa(\bm A)\). The last two lines measure that residual and subtract the
same constant from every atom. That uniform shift makes the pre-clamp sum equal
to \(Q\), up to rounding, and preserves every charge difference \(q_i-q_j\). It
does not re-solve the stationarity equations.

Stored charges then pass through the output clamp. If
\(|q_i|>\code{eeq\_charge\_clamp}\), with the configured value
\(\SI{2.0}{\elementarycharge}\), the stored atom value is set to
\(\pm\SI{2.0}{\elementarycharge}\) and a geometry-choice record is written. A
charge with \(|q_i|=\code{eeq\_charge\_clamp}\) is stored unchanged. The clamp
fires after the uniform charge correction, so the stored sum can move away from
\(Q\). The reported charge statistics are computed from the post-clamp atom
values.

\begin{codebox}
for (size_t i = 0; i < N; ++i) {
    auto& atom = conf.MutableAtomAt(i);
    atom.eeq_cn = cn(i);

    double qi = charges(i);
    if (std::abs(qi) > charge_clamp) {
        qi = (qi > 0) ? charge_clamp : -charge_clamp;
        ++n_clamped;
    }
    atom.eeq_charge = qi;
}
\end{codebox}

\section{Output shape}

EEQ emits scalars, not tensors. It does not call
\code{SphericalTensor::Decompose}, and it does not populate the shared
\([T_0,T_1,T_2]\) tensor layout. A partial charge is a scalar attached to an
atom, but it is not the \(T_0\) part of a shielding tensor.

\begin{codebox}
eeq_charges.npy  // (N,) float64, atom.eeq_charge in elementary charges
eeq_cn.npy       // (N,) float64, atom.eeq_cn coordination number
\end{codebox}

For \(N=0\), \code{Compute} returns \code{nullptr} before any matrix is
allocated. For \(N>0\), \code{WriteFeatures} on a successful result calls
\code{NpyWriter::WriteFloat64} for both arrays with one entry per atom.

\section{Glossary}

\textbf{Smooth coordination number.} A tally with a soft rim --- it fades a
neighbour out across a band of distance instead of clicking from one to zero at
a fixed radius (a neighbour count with a smooth pair weight before the cutoff:
close atoms count nearly one, distant counted atoms count nearly zero, and the
change between them is continuous). In this code, a
pair at
\(R_{ij}=r_{\mathrm{cov},i}+r_{\mathrm{cov},j}\) contributes exactly \(1/2\) to
both atoms before the total is used in \(\chi_i+\kappa_i\sqrt{\mathrm{CN}_i+
10^{-14}}\). Formally, it is the pair sum
\(\mathrm{CN}_i=\sum_{j\ne i}c_{ij}\) with
\(c_{ij}=\frac12\operatorname{erfc}(k_{\mathrm{CN}}
(R_{ij}/(r_{\mathrm{cov},i}+r_{\mathrm{cov},j})-1))\) for counted pairs.

\textbf{Ohno--Klopman kernel.} The \(1/R\) Coulomb curve with its spike at
contact filed down to a finite plateau (regularized to a finite zero-distance
value, the plateau height set by the two chemical hardness parameters). At zero
separation it is \(\sqrt{\eta_i\eta_j}\), while at large separation it tends to
\(1/R_{ij}\); the EEQ matrix writes this same value into both symmetric
off-diagonal slots. Formally,
\(\gamma_{ij}(R)=\big(R^2+1/(\eta_i\eta_j)\big)^{-1/2}\).

\textbf{Cholesky factorization.} A symmetric positive-definite matrix split into
a triangle times its own transpose (a lower triangular \(\bm L\) with
\(\bm A=\bm L\bm L^{\mathsf T}\)), so one dense factorization leaves a pair of
triangular solves for each right-hand side. Here
\code{LLT} forms \(\bm A=\bm L\bm L^{\mathsf T}\), then the two right-hand
sides \(\bm\chi^{\mathrm{eff}}\) and \(\bm1\) reuse that same \(\bm L\).
Formally, Cholesky applies to symmetric positive-definite \(\bm A\) and solves
\(\bm A\bm x=\bm b\) by \(\bm L\bm y=\bm b\) followed by
\(\bm L^{\mathsf T}\bm x=\bm y\).

\textbf{Bordered-system reduction.} A square block with one extra row and column
framing it as a border (an augmented linear system whose extra row and column
enforce a scalar constraint). Rather than factor the whole
bordered matrix, the code solves the inner block twice and recovers the
multiplier with one scalar division. The EEQ code never factors the
\((N+1)\times(N+1)\) KKT matrix; it solves for
\(\bm u=\bm A^{-1}\bm\chi^{\mathrm{eff}}\) and
\(\bm v=\bm A^{-1}\bm1\), then combines them with one scalar \(\lambda\).
Formally, the reduction is
\(\lambda=-(Q+\bm1^{\mathsf T}\bm u)/(\bm1^{\mathsf T}\bm v)\) and
\(\bm q=-(\bm u+\lambda\bm v)\).

\textbf{Conditioning.} The amount a solve magnifies input and rounding errors
on the way from \(\bm b\) to \(\bm x\). In this calculator an ill-conditioned \(\bm A\) makes
\(\bm u\), \(\bm v\), and the scalar sum of \(\bm q\) more sensitive to
roundoff; the code measures only the final charge-sum residual and removes it
by a uniform shift. Formally, the relevant scale is the condition number of
\(\bm A\), which bounds the relative sensitivity of
\(\bm A\bm x=\bm b\) to perturbations in \(\bm A\) and \(\bm b\).

\end{document}
